\documentclass[11pt]{article}
\usepackage[margin=1in]{geometry}
\usepackage{amsmath,amssymb}
\usepackage[hidelinks]{hyperref}

\title{Literature Status: Banach Spaces Without Smooth K-Maps}
\author{FA/Banach run \texttt{fa\_banach\_001}}
\date{2026-07-18}

\begin{document}
\maketitle

\section*{Status}

This is a literature-already-answered status packet, not an original proof
from the run.  Belitskii--Rayskin, arXiv:1703.06299, ask in Section 5,
Open questions, item 1, whether Banach spaces without \(C^\infty\) K-maps
exist.  The existence subquestion has an affirmative later-literature
answer.

\section*{Original Problem}

The original source is:
\begin{quote}
G. Belitskii and V. Rayskin,
\emph{Extension of local smooth maps of Banach spaces},
arXiv:1703.06299.
\end{quote}

The paper defines a \(C^q\)-K-map on a Banach space \(X\) as a globally
defined bounded \(C^q\) representative of the identity germ at the origin.
Equivalently, it is a bounded \(C^q\) map \(H:X\to X\) which agrees with
the identity on a neighborhood of \(0\).  Its open-questions section asks:
for which spaces do \(C^\infty\)-K-maps exist, and do Banach spaces without
\(C^\infty\)-K-maps exist?

\section*{Later Answer}

The supporting source is:
\begin{quote}
G. Belitskii and V. Rayskin,
\emph{A New Method of Extension of Local Maps of Banach Spaces.
Applications and Examples},
arXiv:1706.08513.
\end{quote}

This later paper uses the name \emph{blid map}, defined as a globally
bounded local identity.  Thus the notion is the same as the earlier
K-map notion.  In its section ``More examples and open questions'', it says
that the question whether Banach spaces without smooth blid maps exist has
an affirmative answer, citing D'Alessandro--Hajek and Hajek--Johanis.  The
stated reason is that those works provide Banach spaces which do not allow
\(C^2\)-extension, and hence cannot have a \(C^2\)-blid map.  Therefore
these spaces cannot have a \(C^\infty\)-K-map either.

\section*{Supporting Bibliography}

\begin{itemize}
\item S. D'Alessandro and P. Hajek,
\emph{Polynomial algebras and smooth functions in Banach spaces},
Journal of Functional Analysis \textbf{266} (2014), 1627--1646,
DOI \href{https://doi.org/10.1016/j.jfa.2013.11.017}{10.1016/j.jfa.2013.11.017}.
\item P. Hajek and M. Johanis,
\emph{Smooth Analysis in Banach Spaces},
de Gruyter, 2014.
\end{itemize}

\section*{Verification Notes}

The answer is not contained in arXiv:1703.06299; there the question appears
as an open question.  The later arXiv:1706.08513 source is distinct and
explicitly records the affirmative answer to the same existence question
under the equivalent ``blid map'' terminology.  The implication from no
\(C^2\)-blid map to no \(C^\infty\)-K-map is immediate, since every
\(C^\infty\) K-map is a \(C^2\) bounded local identity.

\section*{Search Bounds}

The local run registry, ledger, solution folders, and attempts were searched
for arXiv:1703.06299, arXiv:1706.08513, K-map, smooth K-map, blid map,
bounded local identity, and the non-even \(\ell_p\) subquestion.  Existing
nearby records cover the harder non-even \(\ell_p\) blid question and a
weak-topology partial result, but no packet for this exact existence-status
subquestion was found.  The later arXiv:1706.08513 source and bibliography
were inspected, and the D'Alessandro--Hajek DOI was checked by web search.

\section*{Limitations}

This packet records only the affirmative answer to the existence of Banach
spaces without smooth K-maps.  It does not classify spaces admitting K/blid
maps.  It does not answer the non-even \(\ell_p\) question, the sphere
question in \(C[0,1]\), the arbitrary-subspace question for \(C[0,1]\), or
the infinite-dimensional Whitney-extension question.

\section*{Human Review Recommendation}

Check the terminology identification between K-maps and blid maps, and
decide whether to attach the underlying D'Alessandro--Hajek JFA PDF if
publisher access permits.  The packet's scope is intentionally narrow:
existence of spaces without smooth K-maps is already answered; the surrounding
classification questions remain open here.

\end{document}
